% Section 6, rewrite of October 22, 2026 (his words 04:12Z): closer to the slides. Two exercises compared, matched supply and
% mortgage credit (quantum and tempo); then who uses the credit (the decomposition); then credit for parents; then cash,
% insurance and a rebated property tax. The +10% supply arm is dropped from the main text. One elasticity (0.63) in the main
% text; eta 1.5 and the slow stock in Appendix D. Main text = free entry (his decision Oct 21 22:36Z).
% NUMBERS ARE PLACEHOLDERS (his word 03:52Z): the slides' runs, OLD SETUP = instant supply eta 0.63, old accounting, 7.9
% calibration, constant-decay state: 20261021_paper_results/p6_values_oldsetup.json and the three fig_p6_*_oldsetup_eta063.pdf
% (p6_figs_oldsetup.py, from 20261017_accounting/transitions/data/prod_*); 2023 tenure of ages 22-29:
% transitions_2023/decay_ship_20261014/tenure_size_22_29_2023_decay.json. \oldv marks every placeholder.
% The 2007 stationary credit figure (\FCC) is at the current calibration. Earlier versions in git.
% Oct 22 15:30Z: FREE-ENTRY numbers filled where the arm has converged (credit, parents, cash, insurance; base TFR):
% 20261021_paper_results/p6_values_fi.json and fig_p6_{paths,transfers}_fi_eta063.pdf (p6_figs_fi.py, from
% 20261017_accounting/transitions/data_fi_torch/P063fi_*.json = Torch pol_063fi/out/runs/*/result.json). Credit is UNDER CHECK
% (1.35% vs about 1.9% on the slides). Still \oldv (red, slides' runs): matched supply, property tax, the 22-29 tenure split, and
% the quantum-tempo figure (needs matched supply).
\providecommand{\pend}{\textcolor{red}{[pending]}}
\providecommand{\oldv}[1]{\textcolor{red}{#1}}
\providecommand{\chk}{\textcolor{red}{$^{\ddagger}$}}
\newcommand{\PRF}{../../output/model/experiments/housing_supply/20261021_paper_results}

\section{Counterfactuals}
\label{sec:policy}
\label{sec:policy-design}

To illustrate the forces at work in the model, I compare two experiments from
2023. In the first, mortgage credit becomes easier: the loan-to-value ratio
rises from $\phi=0.80$ to 1 for all households, so buyers need no down
payment. In the second, the housing supply expands: the supply curve shifts
out by \oldv{6.3\%}, $H_0\to\oldv{1.063}\,H_0$, the size at which births in
2023 rise as much as under credit.\footnote{In the supply experiments the
inherited 2023 distribution is projected onto the states feasible under the
shifted supply; the projected mass is about $10^{-6}$.} The two experiments
are not policy proposals, and the size of the supply expansion is set by the
effect of credit, not by any proposed policy. They lower different costs of
family space: credit lowers the down payment needed to buy a home, and supply
lowers the price of space and with it the per-period cost of a room.
Comparing them at the same effect on impact shows which cost matters for
fertility.

I show that the two experiments raise births by the same amount on impact but
differ in kind: credit moves births earlier, while supply raises the number of
children. Credit moves births earlier because the young use it to own the
small homes they would otherwise rent; it lowers the down payment, not the
per-period cost of the space a child needs. I then extend the comparison to
credit tied to children, which lowers that cost, and to transfers.

\problem{Numbers in red are placeholders from the runs in the slides
(instant supply, earlier calibration and accounting), to be replaced as each
free-entry run lands: the housing supply experiment, the property tax and the split of new
owners by home size. Mortgage credit is under check: it raises births by
1.35\% on impact here, against about 1.9\% in the slides' runs.}

Both experiments start from the 2023 state of the transition in
Section~\ref{sec:quant-transition}: the same households, with the same
wealth, homes and children. In 2023 the change is announced unexpectedly and
believed to be permanent. Households have perfect foresight of the paths of
the house price, rents, the pension and the rebate, which clear their markets
at every date. I compare each experiment with the path the economy follows from
the same state without the change, so that the difference isolates it from
the ongoing transition. By 2023 the preference for children has nearly
completed its decline, and it stays close to its 2023 value along both paths.

Developers enter freely, so the housing stock is on the supply curve at every
date: $H_t=\max\{(1-\delta_H)H_{t-1},\,H_0P_t^{\eta}\}$. The stock follows
the curve $H_0P_t^{\eta}$ unless the curve falls by more than depreciation in
one period; housing cannot be destroyed, so then nothing is built. The stock
responds to the house price with an elasticity $\eta=0.63$
\citep{baum-snowMicrogeographyHousingSupply2024}.
Section~\ref{sec:policy-robust} reports the results with a more elastic
supply and with a stock that adjusts slowly.

I report births in the four-year periods starting in 2023 and 2063 and in the
long run, when the economy reaches its new stationary equilibrium, and the
completed fertility of the cohort aged 18 to 21 in 2023: the children ever
born to these households over their lives. All are relative to the no-policy
path. Table~\ref{tab:policy-main} summarizes the results.

\begin{table}[H]
\centering
\caption{Counterfactuals from 2023}
\label{tab:policy-main}
\small
\begin{tabularx}{\textwidth}{@{}Xrrrcrr@{}}
\toprule
 & \multicolumn{3}{c}{Births} & Completed fertility & \multicolumn{2}{c}{House price} \\
\cmidrule(lr){2-4}\cmidrule(l){6-7}
 & 2023 & 2063 & Long run & 2023 cohort & 2023 & Long run \\
\midrule
\multicolumn{7}{@{}l}{\emph{A. The two experiments}}\\
Mortgage credit\chk & $+1.3$ & $+0.6$ & $+3.0$ & $-1.1$ & $-0.1$ & $+3.2$ \\
Housing supply, \oldv{$+6.3\%$} & \oldv{$+2.0$} & \oldv{$+4.6$} & \oldv{$+6.3$} & \oldv{$+3.1$} & \oldv{$-5.8$} & \oldv{$0.0$} \\
\addlinespace
\multicolumn{7}{@{}l}{\emph{B. Extensions}}\\
Credit for parents & $+5.1$ & $+16.6$ & $+36.1$ & $+4.9$ & $+0.3$ & $+33.5$ \\
Cash transfer & $+8.2$ & $+14.3$ & $+22.8$ & $+7.8$ & $+2.2$ & $+24.0$ \\
Earnings insurance & $+15.2$ & $+27.5$ & $+48.9$ & $+14.0$ & $+3.9$ & $+48.3$ \\
Property tax, rebated & \oldv{$-0.2$} & \oldv{$+1.1$} & \oldv{$+0.5$} & \oldv{$+0.8$} & \oldv{$-6.2$} & \oldv{$-5.0$} \\
\bottomrule
\end{tabularx}
\par\smallskip
\begin{minipage}{0.95\linewidth}
\footnotesize\textit{Notes:} \% differences from the no-policy path from the
same 2023 state, with free entry and $\eta=0.63$. Births: births in the
four-year period starting in the year shown; long run: the new stationary
equilibrium, where births per household are at replacement, so births and the
number of households rise by the same percentage. Completed fertility:
children ever born to the cohort aged 18--21 in 2023. Without policy, period
fertility is 1.63 in 2023 and 1.73 in 2063. Housing supply: the supply curve
shifted out so that births in 2023 rise as much as under mortgage credit. The
extensions are defined in Sections~\ref{sec:policy-parents}
and~\ref{sec:policy-transfers}. \chk\,Under check. Red: placeholders from the
runs in the slides.
\end{minipage}
\end{table}

\subsection{Same impact, different long run}
\label{sec:policy-supply}

Demographers distinguish how many children are born, the \emph{quantum}, from
when they are born, the \emph{tempo} \citep{bongaartsFeeneyQuantumTempo1998}.
Credit is mostly tempo, and supply is quantum. Figure~\ref{fig:policy-qt}
shows the two paths. Both raise births by 1.3\% in 2023. Under credit the
gain is gone within four years, and births fall below the no-policy path, by
0.6\% in 2035, as the births brought forward are missing later. Under supply
births rise in every period, by \oldv{4.6\%} by 2063 and by \oldv{6.3\%} in
the long run.

Supply raises births over time for two reasons. Cohorts that reach
childbearing age after 2023 spend their whole fertile life with cheaper space,
and the additional children form households of their own a generation later.
Credit adds one pulse of births, and its echo fades. Supply adds a steady flow,
and its echoes build up. In the long run the price returns to the level at
which fertility is at replacement, and the number of households rises by the
full \oldv{6.3\%}: the additional rooms are filled by additional families.
Credit also has a small lasting effect, 3.1\% more households in the long
run, with a price 3.2\% higher.

The path of births under credit has four phases (Figure~\ref{fig:policy-qt},
panel a). On impact, credit reaches every renter waiting to buy, at every
fertile age, and some of them have their child now. From 2031 births fall
below the no-policy path: households that had a child early have fewer later,
and those who enter after 2023 only move their births earlier. Births then
rise again. \oldv{About a third} of the rise is the echo, as the children born
in 2023 become young parents, and \oldv{two thirds} comes from new entrants,
who gain more from credit as housing becomes cheaper along the transition.
Credit adds a little to lifetime fertility in the long run because space then
becomes cheap: without policy, the price falls by 45\% from 2023 to the new
stationary equilibrium. At the long-run price, removing the down payment raises births
\oldv{1 percentage point} above replacement; at prices \oldv{about 50\%}
higher, by only \oldv{0.1 to 0.2} points. When space is expensive, the
per-period cost of a family room dominates the decision, and credit, which
changes only the one-time down payment, matters little. Letting renters rent
homes as large as 10 rooms removes \oldv{two thirds} of this long-run effect.

The long run follows from one condition. In a stationary equilibrium the
population is constant only if births per household are at replacement, and
fertility falls with the cost of space. So the long-run price is the one at
which households choose replacement fertility, and the number of households is
the one that fills the housing stock at that price. A policy that raises
fertility at a given price therefore ends up as a higher price and more
households, not as more births per household. A supply expansion leaves the
replacement price unchanged: a 10\% expansion brings the price back to its
no-policy level and leaves 10.0\% more households. Credit for parents raises
fertility at any price, so the price must rise, by 33\%, until fertility is
back at replacement, with 36\% more households.

\begin{figure}[!tb]
\centering
\includegraphics[width=\linewidth]{\PRF/fig_p6_qt_oldsetup_eta063.pdf}
\caption{Mortgage credit and housing supply}
\label{fig:policy-qt}
\par\smallskip
\begin{minipage}{0.92\linewidth}
\footnotesize\textit{Notes:} \% differences from the no-policy path from the
same 2023 state, births in each four-year period (panel a) and the number of
households (panel b), to the new stationary equilibrium. Mortgage credit: LTV
80 to 100 for all households. Housing supply: the supply curve shifted out by
\oldv{6.3\%}. \oldv{Placeholder figure: runs in the slides, until the housing
supply run lands.}
\end{minipage}
\end{figure}

Table~\ref{tab:policy-tempo} follows the cohort aged 18 to 21 in 2023, which
separates the number of children from their timing. Credit lowers its
completed fertility by 1.1\% and the share with a first birth by 1.3\%, and
brings first births forward by 0.08 years. The supply expansion
raises completed fertility by \oldv{3.1\%} and the share with a first birth by
\oldv{2.4\%}. Supply also brings first births forward, by \oldv{0.14} years: a
cheaper room lowers the per-period cost of the commitment, so households
commit sooner as well as more often.

\begin{table}[H]
\centering
\caption{The cohort aged 18--21 in 2023}
\label{tab:policy-tempo}
\small
\begin{tabularx}{\textwidth}{@{}Xrrr@{}}
\toprule
 & Completed fertility & Share with a first birth & Mean age at first birth \\
\midrule
No policy & 1.640 & 0.673 & 27.15 \\
Mortgage credit\chk & 1.622 ($-1.1\%$) & 0.665 ($-1.3\%$) & 27.07 \\
Housing supply & \oldv{1.693} ($\oldv{+3.1\%}$) & \oldv{0.693} ($\oldv{+2.4\%}$) & \oldv{26.95} \\
Credit for parents & 1.720 ($+4.9\%$) & 0.701 ($+4.1\%$) & 26.88 \\
\bottomrule
\end{tabularx}
\par\smallskip
\begin{minipage}{0.95\linewidth}
\footnotesize\textit{Notes:} Children ever born and the share with a first
birth, per household of the cohort, summed over ages 18 to 45. Mean age at
first birth: first-birth rates weighted by the midpoints of the four-year age
cells. \chk\,Under check. Red: housing supply from the runs in the slides
(a different no-policy path: 1.643, 0.676, 27.09).
\end{minipage}
\end{table}

\subsection{Who uses the credit}
\label{sec:policy-credit}

Why does credit move births earlier rather than raise their number? Credit
lowers the equity needed to buy a home, $(1-\phi)P_th^O$, not the per-period
cost of the space a child needs. Rent equals the user cost of owning
\eqref{eq:m-rent}, so a family can rent that space at about the per-period
cost of owning it. By \eqref{eq:m-credit}, credit raises the probability of
trying for a child only if it is worth more to a household with the additional
child than without it. It is worth more without the child when it mainly helps
households without children to buy.

The households that use the credit confirm this. Among households aged 22 to
29 in 2023, ownership rises from \oldv{49\%} to \oldv{73\%}. \oldv{Seven in eight} of the
new owners buy a home of 2 or 4 rooms. The share owning a home of 6 rooms or
more, the size families choose, rises by only \oldv{3} percentage points.
Across all households, ownership rises from 74.9\% to 82.8\% in 2023, and the
mean size of an owned home falls from 6.05 to 5.73 rooms. Total demand for
space falls by 0.09\%, and the price with it, by 0.14\%. The price barely
matters for this result: at the no-policy prices, credit raises births in
2023 by 1.38\%, against 1.35\% in equilibrium.

Credit buys small homes early. Among households aged 22 to 25 in 2027,
ownership rises from \oldv{46\%} to \oldv{68\%}, and the share owning a
two-room home from \oldv{8\%} to \oldv{25\%}, while births per 100 households
at those ages stay at \oldv{22.2}. A small home accommodates a young family,
but the households that buy one early are not those whose fertility was held
back. The pattern persists along the transition. In 2059 credit raises
ownership at ages 22 to 37 by \oldv{21} percentage points, and young owners
live in smaller homes, \oldv{5.64} rooms against \oldv{6.36}, with more debt,
\oldv{0.68} of the value of the home against \oldv{0.55}. A supply expansion
instead gives young owners larger homes, \oldv{6.65} rooms, with the same debt.
Credit buys earlier, smaller and more leveraged ownership; supply buys more
space.

The PSID offers a partial check on this mechanism. Among childless first-time
buyers aged 22 to 33, those who bought a home of four rooms or fewer were 14
percentage points less likely to have a first birth within six years
(standard error 6.7). The gap goes in the model's direction, but without an
instrument it is as consistent with selection, couples not planning children
buying small, as with an effect of the home.

To isolate this margin from the response of prices, I raise $\phi$ to 1 in the
2007 stationary equilibrium, at the same price. Ownership at ages 25 to 34
rises from 58\% to 85\%. Figure~\ref{fig:policy-credit-own} shows that the new
owners are young and buy small homes: at ages 22 to 25, the share owning a
home of 2 or 4 rooms rises from 12\% to 49\% (panel a). The share owning a
home of 6 rooms or more rises by one percentage point among the youngest and
falls slightly at older ages (panel b). Births move slightly earlier, and
completed fertility is unchanged (panel c). The young use the credit to own
the small homes they would otherwise have rented. Credit changes who owns and
when, not the cost of a child's space.

\begin{figure}[!tb]
\centering
\includegraphics[width=\linewidth]{\FCC/fig_s5_credit.pdf}
\caption{Mortgage credit, ownership and births}
\label{fig:policy-credit-own}
\par\smallskip
\begin{minipage}{0.92\linewidth}
\footnotesize\textit{Notes:} 2007 stationary equilibrium at a loan-to-value
ratio of 80\% and of 100\%, same price. Share of households owning a home of
2 or 4 rooms (panel a) and of 6 rooms or more (panel b), and births per
household in each four-year age group (panel c), by age of the head.
\end{minipage}
\end{figure}

\subsection{Credit for parents}
\label{sec:policy-parents}

If credit fails because it reaches households whose fertility was not held
back, credit tied to children should do better. I raise the loan-to-value
ratio to 1 only for households with a child at home. This lowers the cost of
the child rather than the cost of waiting. Births in 2023 rise by 5.1\%,
several times as much as under credit for all households, at nearly the same
price, 0.3\% higher (Figure~\ref{fig:policy-paths}). The completed
fertility of the 2023 cohort rises by 4.9\%, and the share with a first
birth by 4.1\% (Table~\ref{tab:policy-tempo}). Births keep rising, by
17\% by 2063. As the larger cohorts demand more space, the price rises, by
3.5\% by 2063 and 33\% in the long run, when there are 36\% more
households.

The policy works through two channels, which the experiment does not
separate. First, a birth unlocks ownership without a down payment, so the
policy rewards the birth itself. Second, a parent who owns can borrow up to
the full value of the home in every period with a child at home, a buffer in
the state that carries the commitment.

\begin{figure}[!tb]
\centering
\includegraphics[width=0.9\linewidth]{\PRF/fig_p6_paths_fi_eta063.pdf}
\caption{Mortgage credit, housing supply and credit for parents}
\label{fig:policy-paths}
\par\smallskip
\begin{minipage}{0.92\linewidth}
\footnotesize\textit{Notes:} \% differences from the no-policy path from the
same 2023 state. Births in each four-year period (panel a), completed
fertility of the cohort aged 18--21 in the year shown (panel b), number of
households (panel c) and house price (panel d). Credit for parents: LTV 100
for households with a child at home. \oldv{Housing supply: pending. Mortgage
credit: under check.}
\end{minipage}
\end{figure}

\subsection{Cash, insurance and a property tax}
\label{sec:policy-transfers}

I compare the housing policies with three policies that work through
households' resources rather than through housing. \emph{Earnings insurance}
pays every working household a quarter of the shortfall of its earnings below
the median earnings state at its age. \emph{Cash transfer} pays the expected
value of that insurance at each age to every working household. Neither is
financed within the model, so neither is a comparison at equal fiscal cost.
\emph{Property tax}: the annual tax rises from 1.06\% to 2\% of value, and the
revenue is rebated equally to all households.

Transfers raise births far more than any of the housing policies
(Figure~\ref{fig:policy-transfers}). Cash raises births in 2023 by
8.2\% and the completed fertility of the 2023 cohort by 7.8\%;
insurance by 15.2\% and 14.0\%. The difference, 6 to 7
percentage points, is the value of insuring the downside of earnings rather
than paying the same amount to everyone: a child is a long commitment, and
insurance makes it safer. Both raise the price, by 24\% and 48\% in
the long run, as the larger population demands more space.

Insurance is a partial version of lower earnings risk, which gives the
benchmark. Lowering the standard deviation of the four-year earnings shock
from 0.48 to 0.35 from 2023 on, with mean earnings unchanged, raises births in
2023 by \oldv{11.3\%} and the completed fertility of the 2023 cohort by
\oldv{16.0\%}, from \oldv{1.64} to \oldv{1.91}. The additional demand for
family space raises the price, by \oldv{11\%} by 2063 and \oldv{49\%} in the
long run.

The property tax lowers the price by \oldv{6.2\%} in 2023, close to the supply
expansion, but births do not move (\oldv{$-0.2\%$}) and the completed
fertility of the 2023 cohort rises by only \oldv{0.8\%}. The tax lowers the price
because it raises the cost of holding housing, so the per-period cost of a room
rises: rents are \oldv{13.5\%} higher in 2023. The rebate returns the revenue per
household, not per room, so a family that needs more space pays more. A lower
price is not cheaper space.

\begin{figure}[!tb]
\centering
\includegraphics[width=\linewidth]{\PRF/fig_p6_transfers_fi_eta063.pdf}
\caption{Cash, insurance and a property tax}
\label{fig:policy-transfers}
\par\smallskip
\begin{minipage}{0.92\linewidth}
\footnotesize\textit{Notes:} \% differences from the no-policy path from the
same 2023 state, births in each four-year period (panel a) and house price
(panel b). \oldv{Property tax: pending.}
\end{minipage}
\end{figure}

\subsection{A more elastic supply and a slowly adjusting stock}
\label{sec:policy-robust}

Appendix~\ref{app:policy-eta} repeats the experiments with $\eta=1.5$. With a
more elastic supply, a supply expansion lowers the price less, and the
additional demand from credit for parents raises it less. \pend{}

Appendix~\ref{app:policy-slow} reports the experiments when the stock adjusts
slowly, as in \citet{kaplanMitmanViolanteHousing2020}. Developers then need
permits, which are in fixed supply, so only new construction responds to the
price within a period, and the stock responds with an elasticity of about
$\delta_H\eta$. Mortgage credit raises births by 1.7\% in 2023 and leaves the
completed fertility of the 2023 cohort about unchanged ($-0.2\%$), and credit
for parents raises it by 5.9\%. A supply expansion works more slowly than with
free entry: a 10\% expansion lowers the price by 3.0\% in 2023 and raises
births by 0.4\%, and births are 3.2\% higher by 2063, as construction keeps
adding space.
